% FORMALIZACION_REUNIDA. Desarrollo ampliado de las tres composiciones de
% EXPLORACION_AMBIVALENCIA_ESPEJO_DIAGONAL_20260918, con dependencias de X.
% Las coordenadas pi, phi, e, alpha y el registro K son salidas generadas.
% Este capítulo define lectores posteriores y no altera el generador TPK.
\chapter{Nonadic holonomy, reciprocity and recovery of geometric readers}
\label{hrr:capitulo}

The nonadic connection, areal--volumetric incidence and radial reciprocity enter different realizations of the same enriched state. The connection preserves the oriented sequence of nine phases and extends memory on return. Incidence records transitions between two reading modes. Radial reciprocity exchanges the two orientations of a geometric form of fixed area. Their composition is expressed through three precise results: invariance of a mark quotient at every depth, factorization of a determinant-one deformation and reconstruction of anisotropy from a normalized length and its orientation.

The construction order remains
\[
 \begin{gathered}
 \mathrm{APP}\longrightarrow\mathrm{TRIT}\longrightarrow\mathrm{TPK}
 \longrightarrow\text{enriched state},\\
 \text{enriched state}
 \longrightarrow\text{joint discrete structure of the continuum},\\
 \text{joint discrete structure of the continuum}
 \longrightarrow\text{readers composed here}.
 \end{gathered}
\]
The coordinates \(\pi,\varphi,e,\alpha\) and register \(K\) appearing in the formulas are outputs of the preceding regional generation. The tritic calendar already determines consecutive pairs of additive and multiplicative modes. The language of matrices, hyperbolic functions and exterior minors then makes explicit the realizations of that generated structure.

\section{Modal incidence and regional marks on finite returns}
\label{hrr:incidencia}

\subsection{Enumeration of returns between sheets}

The temporal block \((+,-,0)\) induces the modal sequence \((\Sigma,\Pi,\Pi)\). Its consecutive pairs, including the return link, are \((\Sigma,\Pi)\), \((\Pi,\Pi)\) and \((\Pi,\Sigma)\). The primitive incidence, in areal--volumetric order, is therefore
\begin{equation}
 F_{\mathrm{av}}=
 \begin{pmatrix}0&1\\1&1\end{pmatrix}.
 \label{hrr:incidencia-matriz}
\end{equation}
This matrix records the calendar pairs. The memory-one symbolic envelope it determines preserves that incidence; additional phase, orientation and register restrictions remain in the TPK state. The envelope count and survival of a complete enriched history are thus distinguished.

Let \(f_0=0\), \(f_1=1\) and \(f_{n+1}=f_n+f_{n-1}\) for \(n\geq1\). These integers express here the counts produced by powers of the incidence. Conventional identification of this sequence remains subsequent to construction of the calendar.

\begin{proposition}[Powers of incidence and diagonal returns]
\label{hrr:potencias}
For every \(n\geq1\),
\begin{equation}
 F_{\mathrm{av}}^n=
 \begin{pmatrix}f_{n-1}&f_n\\f_n&f_{n+1}\end{pmatrix}.
 \label{hrr:potencias-formula}
\end{equation}
In particular, returns of length \(n\) to the areal and volumetric sheets have multiplicities \(f_{n-1}\) and \(f_{n+1}\), respectively.
\end{proposition}
\begin{proof}
For \(n=1\), the formula is the definition of \(F_{\mathrm{av}}\), since \(f_2=1\). Assuming validity for \(n\), right multiplication gives
\[
 \begin{pmatrix}f_{n-1}&f_n\\f_n&f_{n+1}\end{pmatrix}
 \begin{pmatrix}0&1\\1&1\end{pmatrix}
 =\begin{pmatrix}f_n&f_{n-1}+f_n\\
 f_{n+1}&f_n+f_{n+1}\end{pmatrix}
 =\begin{pmatrix}f_n&f_{n+1}\\f_{n+1}&f_{n+2}\end{pmatrix}.
\]
Induction proves the identity. For an incidence matrix with entries zero or one, the multiplication rule sums sequences of compatible transitions between two positions. By induction on length, entry \((i,j)\) of the power counts sequences of length \(n\) from \(i\) to \(j\). The two diagonal entries count the stated returns.
\end{proof}

\subsection{Regional marking of the return support}

Let \(e_{\mathrm{ar}}=(1,0)^{\mathsf T}\) and \(e_{\mathrm{vol}}=(0,1)^{\mathsf T}\). For a positive coordinate \(c\), produced before this evaluation, define the mark reader
\[
 D_c=c\,e_{\mathrm{ar}}e_{\mathrm{ar}}^{\mathsf T}
       +e_{\mathrm{vol}}e_{\mathrm{vol}}^{\mathsf T}
     =\operatorname{diag}(c,1).
\]
Diagonality preserves both sheets; the volumetric entry expresses its unit normalization. The areal mark \(c\) is applied once, after the return count. Its evaluation is
\begin{equation}
 \Theta_{c,n}=
 \frac{\langle e_{\mathrm{ar}},D_c F_{\mathrm{av}}^n
 e_{\mathrm{ar}}\rangle}
 {\langle e_{\mathrm{vol}},F_{\mathrm{av}}^n
 e_{\mathrm{vol}}\rangle}
 =c\,\frac{f_{n-1}}{f_{n+1}},\qquad n\geq1.
 \label{hrr:marca}
\end{equation}
The choice \(c=\pi\) coincides with the boundary reader of the areal--volumetric return. The choice \(c=e\) defines the complementary subsequent reader considered in this chapter. Both use the same return support and the regional coordinate already generated.

\begin{theorem}[Depth-independent mark quotient]
\label{hrr:cociente} For any positive marks \(c,d\) and every \(n\geq2\),
\begin{equation}
 \frac{\Theta_{c,n}}{\Theta_{d,n}}=\frac cd.
 \label{hrr:cociente-exacto}
\end{equation}
In particular,
\begin{equation}
 \frac{\Theta_{\pi,n}}{\Theta_{e,n}}=\frac\pi e,
 \qquad
 \lim_{n\to\infty}\Theta_{\pi,n}=\frac\pi{\varphi^2},
 \qquad
 \lim_{n\to\infty}\Theta_{e,n}=\frac e{\varphi^2}.
 \label{hrr:pi-e-phi}
\end{equation}
\end{theorem}
\begin{proof}
The recurrence implies \(f_j>0\) for \(j\geq1\). For \(n\geq2\), the quotient \(r_n=f_{n-1}/f_{n+1}\) is positive. By \eqref{hrr:marca}, the two evaluations are \(cr_n\) and \(dr_n\), so division cancels exactly the same return factor. The preceding identification of the regional self-scaling coordinate with the stable ratio of this incidence gives \(r_n\to\varphi^{-2}\). Multiplying by the constant marks \(\pi\) and \(e\) yields both limits. The finite quotient, by contrast, already has its exact value before passage to the limit.
\end{proof}

The first length has \(f_0=0\): both evaluations vanish and their quotient lies outside the domain of \eqref{hrr:cociente-exacto}. The restriction \(n\geq2\) expresses a property of the count and avoids extending a simplification to a zero denominator.

\begin{center}
\begin{tabular}{c|ccccc}
\(n\)&2&3&4&5&6\\\hline
\(r_n\)&\(1/2\)&\(1/3\)&\(2/5\)&\(3/8\)&\(5/13\)\\
\(\Theta_{\pi,n}\)&\(\pi/2\)&\(\pi/3\)&\(2\pi/5\)&\(3\pi/8\)&\(5\pi/13\)\\
\(\Theta_{e,n}\)&\(e/2\)&\(e/3\)&\(2e/5\)&\(3e/8\)&\(5e/13\)
\end{tabular}
\end{center}
The table exhibits two simultaneous facts: each reader changes with depth, while their comparison remains fixed. Invariance pertains to a relationship between readers, not to immobility of the state supporting the reading.

\subsection{Exact convergence bound to the stable ratio}

The already established correspondence of the generated coordinate \(\varphi\) with the positive eigenvalue of \(F_{\mathrm{av}}\) gives \(\varphi^2=\varphi+1\) and \(\varphi>1\). The limit can be quantified without introducing a target decimal expansion.

\begin{proposition}[Error of the return quotient]
\label{hrr:error-retorno} Setting \(q=-\varphi^{-2}\), for \(n\geq1\) we have
\begin{align}
 r_n&=\varphi^{-2}\frac{1-q^{n-1}}{1-q^{n+1}},
 \label{hrr:rn-exacto}\\
 r_n-\varphi^{-2}
 &=\varphi^{-2}
 \frac{q^{n-1}(q^2-1)}{1-q^{n+1}}.
 \label{hrr:rn-error}
\end{align}
In particular,
\begin{equation}
 |r_n-\varphi^{-2}|
 \leq\varphi^{-2}
 \frac{|q|^{n-1}(1-q^2)}{1-|q|^{n+1}}.
 \label{hrr:rn-cota}
\end{equation}
\end{proposition}
\begin{proof}
Let \(\psi=-\varphi^{-1}\). From \(\varphi^2=\varphi+1\) follows \(\psi^2=\psi+1\). The sequence \(b_n=(\varphi^n-\psi^n)/(\varphi-\psi)\) satisfies \(b_0=0\), \(b_1=1\) and the same recurrence as \(f_n\). Induction gives \(b_n=f_n\). Dividing the expression for \(f_{n-1}\) by that for \(f_{n+1}\), and using \(\psi/\varphi=q\), yields \eqref{hrr:rn-exacto}. Subtracting the limit produces \eqref{hrr:rn-error}. Finally, \(|q|<1\), \(1-q^2>0\) and \(|1-q^{n+1}|\geq1-|q|^{n+1}>0\) give the bound.
\end{proof}

This formula is an exact reading of the convergence rate of the incidence already generated. The marks receive the error multiplied by \(c\), but their quotient retains the exact cancellation of Theorem~\ref{hrr:cociente}.

\section{Multiplicative composition of the coupling flow}
\label{hrr:flujo}

\subsection{Generated direction and elimination of the uniform component}

The dodecaphase register \(K\), with its marking, produces the direction
\[
 u_K=P_3P_{11}K,\qquad u_K\neq0,\qquad
 \langle u_K,\mathbf1\rangle=0.
\]
The projectors are the operators already constructed in the corresponding realization of the register. This section fixes their result and writes
\[
 v=\frac{u_K}{\|u_K\|},\qquad
 H_K=\operatorname{diag}(v_1,\ldots,v_{12}).
\]
Then \(\sum_i v_i=0\), \(\sum_i v_i^2=1\) and \(\operatorname{tr}H_K=0\). For a positive mark \(r\) and a real parameter \(t\), define
\begin{equation}
 g_{K,r}(t)=\exp\bigl(t\log(r)H_K\bigr)
 =\operatorname{diag}\bigl(r^{tv_1},\ldots,r^{tv_{12}}\bigr).
 \label{hrr:g-definicion}
\end{equation}
The direction comes from the register; the mark determines the logarithmic scale of deformation. Parameter \(t\) parametrizes this geometric flow. Identification with a physical clock requires the corresponding temporal realization.

\begin{theorem}[Mark homomorphism and determinant preservation]
\label{hrr:homomorfismo}
For \(a,b>0\) and \(s,t\in\mathbb R\),
\begin{align}
 g_{K,ab}(t)&=g_{K,a}(t)g_{K,b}(t),
 \label{hrr:producto}\\
 g_{K,a/b}(t)&=g_{K,a}(t)g_{K,b}(t)^{-1},
 \label{hrr:cociente-flujo}\\
 g_{K,r}(s+t)&=g_{K,r}(s)g_{K,r}(t),
 \label{hrr:grupo-t}\\
 \det g_{K,r}(t)&=1.
 \label{hrr:determinante}
\end{align}
Moreover, \(g_{K,r}(t)^{-1}=g_{K,r}(-t)=g_{K,r^{-1}}(t)\).
\end{theorem}
\begin{proof}
All matrices are diagonal. Entry \(i\) of the first equality is \((ab)^{tv_i}=a^{tv_i}b^{tv_i}\), since \(\log(ab)=\log a+\log b\) for positive bases. Every entry is strictly positive, so the matrix is invertible, and substitution \(b\mapsto b^{-1}\) proves the second equality. The third follows from \(r^{(s+t)v_i}=r^{sv_i}r^{tv_i}\). The determinant is the product of the twelve entries:
\[
 \prod_{i=1}^{12}r^{tv_i}
 =r^{t\sum_i v_i}=r^0=1.
\]
The inverse formula again consists of inverting every entry.
\end{proof}

\begin{corollary}[Composition of ambivalence and regional mirror]
\label{hrr:composicion}
For every \(n\geq2\),
\begin{equation}
 g_{K,\Theta_{\pi,n}}(t)
 g_{K,\Theta_{e,n}}(t)^{-1}
 =g_{K,\pi/e}(t).
 \label{hrr:composicion-finita}
\end{equation}
In the stable limit,
\begin{equation}
 g_{K,\pi/\varphi^2}(t)
 =g_{K,\pi/e}(t)\,g_{K,e/\varphi^2}(t).
 \label{hrr:composicion-limite}
\end{equation}
\end{corollary}
\begin{proof}
Both finite marks are positive for \(n\geq2\). Applying \eqref{hrr:cociente-flujo} and \eqref{hrr:cociente-exacto}, their matrix quotient is the flow of mark \(\pi/e\). For the second equality, it suffices to use \((\pi/e)(e/\varphi^2)=\pi/\varphi^2\) in \eqref{hrr:producto}. It also follows as the limit of the finite identity, since the exponential and logarithm are continuous on the domain of positive bases.
\end{proof}

The equality compares deformations with the same direction \(u_K\). Internal mirror exchange of the regional channels retains its own domain on the enriched state; the ratio \(\pi/e\) is its subsequent scalar reading. The holonomy operator \(\Gamma_9\) continues to act on phase and memory. These domains determine the meaning of the composition and preserve the register accompanying return.

\subsection{Transport of minors and exterior invariants}

Let \(X\) be a real matrix of rank \(k\) with twelve columns, and \(\Delta_I(X)\) the column minor specified by a set \(I\) of \(k\) indices. Right action by \(g_{K,r}(t)\) scales column \(i\) by \(r^{tv_i}\).

\begin{proposition}[Signs, strata and quotient invariants]
\label{hrr:plucker}
We have
\begin{equation}
 \Delta_I\bigl(Xg_{K,r}(t)\bigr)
 =r^{t\sum_{i\in I}v_i}\Delta_I(X).
 \label{hrr:plucker-formula}
\end{equation}
The signs and vanishing pattern of all minors are preserved. If \(I,J,P,Q\) satisfy \(\mathbf1_I+\mathbf1_J=\mathbf1_P+\mathbf1_Q\), then, wherever \(\Delta_P(X)\Delta_Q(X)\neq0\), the quotient
\begin{equation}
 \frac{\Delta_I\Delta_J}{\Delta_P\Delta_Q}
 \label{hrr:cociente-menores}
\end{equation}
is invariant under the flow.
\end{proposition}
\begin{proof}
Multilinearity of the determinant multiplies the minor by the product of its column factors, proving \eqref{hrr:plucker-formula}. This product is positive and therefore preserves the minor's sign and vanishing. In \eqref{hrr:cociente-menores}, the total exponent is
\[
 t\sum_{i=1}^{12}
 (\mathbf1_I(i)+\mathbf1_J(i)-\mathbf1_P(i)-\mathbf1_Q(i))v_i=0.
\]
The denominator remains nonzero and cancellation is exact.
\end{proof}

For \(k=2\), a concrete application is
\[
 \frac{\Delta_{12}\Delta_{34}}{\Delta_{13}\Delta_{24}}.
\]
Each of the four indices occurs once in the numerator and once in the denominator. In a configuration where the denominator minors are nonzero, this quantity remains fixed even though each individual minor changes. Exterior conservation thus possesses explicit observables in addition to the global determinant.

\subsection{Lattice lifting and duality}

The marked lattice realization of the preceding chapter uses the orthogonal sum of twelve planes. Write
\[
 G_{K,r}(t)=\bigoplus_{i=1}^{12}r^{tv_i}I_2.
\]
If \(\Lambda\) is the reference self-dual Euclidean lattice already obtained, define \(\Lambda_{r,t}=G_{K,r}(t)\Lambda\).

\begin{proposition}[Covolume and reciprocity of deformed lattices]
\label{hrr:redes} The covolume of \(\Lambda_{r,t}\) equals that of \(\Lambda\), and
\begin{equation}
 \Lambda_{r,t}^{*}=\Lambda_{r,-t}.
 \label{hrr:red-dual}
\end{equation}
The factorization \eqref{hrr:composicion-limite} lifts to these twenty-four-dimensional operators.
\end{proposition}
\begin{proof}
Each radius occurs twice in the determinant. Therefore,
\[
 \det G_{K,r}(t)=\prod_i r^{2tv_i}=1.
\]
The linear change multiplies the covolume by the absolute value of the determinant, hence preserves it. For any invertible operator \(G\), the condition \(y\in(G\Lambda)^*\) is equivalent to \(\langle y,Gx\rangle\in\mathbb Z\) for every \(x\in\Lambda\). It is equivalent to \(G^{\mathsf T}y\in\Lambda^*\), hence \((G\Lambda)^*=G^{-\mathsf T}\Lambda^*\). Here \(\Lambda^*=\Lambda\), \(G\) is positive diagonal and \(G^{-\mathsf T}=G_{K,r}(-t)\), proving \eqref{hrr:red-dual}. Factorization is verified on each two-dimensional block exactly as in \eqref{hrr:producto}.
\end{proof}

This conclusion preserves covolume and relates the deformation to its dual. Integrality of scalar products and evenness of norms are additional lattice properties. Theta--Mellin spectral reciprocity uses precisely the duality proved, together with Gaussian summability and the Poisson formula developed in their own chapter.

\section{Normalized diagonal and restitution of radial orientation}
\label{hrr:diagonal}

\subsection{Separation of scale and shape}

The antecedent radial chart starts from the generated angular coordinates \(A>0\) and \(C^*\), in the domain \(|C^*|<A\). Define
\[
 \xi=\frac{C^*}{A},\qquad
 \eta=\operatorname{artanh}\xi.
\]
The variable \(\eta\) denotes shape anisotropy. Its meaning differs from that of the return mantissa \(\eta_{\mathrm{ret}}\) entering the action section.

Let \(s>0\) be the areal scale produced by that section. In the action-ellipse realization, \(s=2\hbar\). The positive semiaxes are
\begin{equation}
 a=\sqrt{s}\,e^{\eta/2},\qquad
 b=\sqrt{s}\,e^{-\eta/2},\qquad ab=s.
 \label{hrr:semiejes}
\end{equation}
Reciprocity exchanges \(a\) and \(b\). Its radial coordinate is \(R_*=\sqrt{b/a}=e^{-\eta/2}\). The length considered below is the diagonal of the rectangle with sides \(a,b\):
\[
 d=\sqrt{a^2+b^2},\qquad D=\frac{d}{\sqrt{ab}}.
\]
Normalization removes the areal scale and preserves the rectangle's shape up to exchange of its sides.

\begin{theorem}[Diagonal, anisotropy and reciprocity]
\label{hrr:diagonal-teorema} In the preceding domain,
\begin{equation}
 D^2=2\cosh\eta=R_*^2+R_*^{-2},
 \qquad
 D=\frac{\sqrt2}{(1-\xi^2)^{1/4}}.
 \label{hrr:diagonal-formula}
\end{equation}
We have \(D\geq\sqrt2\), and the following are equivalent
\[
 D=\sqrt2,\qquad \eta=0,\qquad \xi=0,\qquad R_*=1,
 \qquad a=b.
\]
For \(D>\sqrt2\), the additional datum \(\varepsilon=\operatorname{sgn}\eta\in\{-1,+1\}\) uniquely determines
\begin{align}
 \eta&=\varepsilon\operatorname{arcosh}(D^2/2),
 \label{hrr:eta-recuperada}\\
 \xi&=\varepsilon\sqrt{1-4/D^4},
 \label{hrr:xi-recuperada}\\
 R_*&=\exp\!\left[-\frac{\varepsilon}{2}
                   \operatorname{arcosh}(D^2/2)\right].
 \label{hrr:radio-recuperado}
\end{align}
\end{theorem}
\begin{proof}
Dividing \(a^2+b^2\) by \(ab=s\) in \eqref{hrr:semiejes} yields
\[
 D^2=e^{\eta}+e^{-\eta}=2\cosh\eta.
\]
Since \(R_*^2=e^{-\eta}\), the sum is also \(R_*^2+R_*^{-2}\). From \(\xi=\tanh\eta\) and \(1-\tanh^2\eta=\cosh^{-2}\eta\), with \(\cosh\eta>0\), follows \(\cosh\eta=(1-\xi^2)^{-1/2}\). The positive root supplies the second formula in \eqref{hrr:diagonal-formula}.

The inequality \(\cosh\eta\geq1\), with equality only for \(\eta=0\), proves the bound and its equivalences through the definitions. For \(D>\sqrt2\), \(\cosh\) determines \(|\eta|=\operatorname{arcosh}(D^2/2)\), and orientation \(\varepsilon\) determines its sign. Moreover,
\[
 \xi^2=1-\cosh^{-2}\eta=1-\frac4{D^4}.
\]
Monotonicity of \(\tanh\) implies \(\operatorname{sgn}\xi=\operatorname{sgn}\eta\), hence \eqref{hrr:xi-recuperada}. Substitution of \eqref{hrr:eta-recuperada} into \(R_*=e^{-\eta/2}\) proves the last formula. Every step is unambiguous in the stated domain.
\end{proof}

\begin{corollary}[Exact recovery of the semiaxes]
\label{hrr:semiejes-recuperacion} The data \((s,D,\varepsilon)\), with \(s>0\) and \(D>\sqrt2\), recover the ordered semiaxes through
\begin{equation}
 a^2=\frac{s}{2}\left(D^2+\varepsilon\sqrt{D^4-4}\right),
 \qquad
 b^2=\frac{s}{2}\left(D^2-\varepsilon\sqrt{D^4-4}\right).
 \label{hrr:semiejes-radicales}
\end{equation}
If \(D=\sqrt2\), recovery is \(a=b=\sqrt{s}\).
\end{corollary}
\begin{proof}
The equations \(a^2+b^2=sD^2\) and \(a^2b^2=s^2\) determine \(a^2,b^2\) as the positive roots of \(z^2-sD^2z+s^2=0\). Its roots are the two quantities \(s(D^2\pm\sqrt{D^4-4})/2\). The sign of \(a-b\) agrees with that of \(\eta\), so \(\varepsilon\) assigns each root to its semiaxis. The roots are positive because \(D^2>\sqrt{D^4-4}\). Equality produces a double root \(s\).
\end{proof}

\subsection{Exact example of two reciprocal realizations}

Take a test shape with \(\xi=3/5\). This example verifies the geometric reader; the concrete assignment of an HMT section is obtained from its own generated coordinates. We compute
\[
 \eta=\operatorname{artanh}(3/5)=\log2,\qquad
 R_*=\frac1{\sqrt2},\qquad
 D^2=2+\frac12=\frac52.
\]
For scale \(s\), the semiaxes are \(a=\sqrt{2s}\) and \(b=\sqrt{s/2}\). The reciprocal shape has \(\xi=-3/5\), \(\eta=-\log2\), \(R_*=\sqrt2\) and exchanged semiaxes. Both have the same normalized diagonal \(D=\sqrt{5/2}\) and the same product \(ab=s\). The orientation datum distinguishes these two realizations exactly.

The diagonal therefore represents anisotropy intensity. The oriented register reconstructs the ordered shape. The areal scale, in turn, recovers its size. This separation specifies which datum each reader preserves and which datum the archive supplies.

\subsection{Poincaré and Klein charts}

The Poincaré chart of the radial diameter and its Klein chart are
\begin{equation}
 w=\frac{R_*^2-1}{R_*^2+1}
   =-\tanh(\eta/2),\qquad
 x_{\mathrm K}=\frac{2w}{1+w^2}=-\tanh\eta=-\xi.
 \label{hrr:cartas}
\end{equation}
The domains \(R_*>0\), \(-1<w<1\) and \(-1<x_{\mathrm K}<1\) correspond through these maps.

\begin{proposition}[Radial reciprocity in both charts]
\label{hrr:reciprocidad-cartas} Inversion \(R_*\mapsto R_*^{-1}\) acts through
\[
 \eta\mapsto-\eta,\qquad w\mapsto-w,
 \qquad x_{\mathrm K}\mapsto-x_{\mathrm K}.
\]
The diagonal and areal scale remain invariant. The pair \((D,\varepsilon)\) recovers both disk coordinates through \eqref{hrr:eta-recuperada} and \eqref{hrr:cartas}.
\end{proposition}
\begin{proof}
The identity \(R_*=e^{-\eta/2}\) transforms inversion into a sign change of \(\eta\). Both hyperbolic functions in \eqref{hrr:cartas} are odd, hence change sign. Formula \(D^2=R_*^2+R_*^{-2}\) is invariant, and the product of the semiaxes remains \(s\). The inverse formulas of Theorem~\ref{hrr:diagonal-teorema} complete recovery.
\end{proof}

In the example \(\xi=3/5\), we obtain \(w=-1/3\) and \(x_{\mathrm K}=-3/5\). The reciprocal shape has \(w=1/3\) and \(x_{\mathrm K}=3/5\). Coincidence of the diagonal thus accompanies two opposite positions on the same diameter.

\subsection{Radial inversion and reflection in the complex chart}

The positive radius also enters the complex reciprocity chart. For \(r>0\) and an argument \(\theta\), let
\[
 z=re^{i\theta},\qquad s(r,\theta)=\frac1{1-z},
 \qquad z\neq1.
\]
The domain excludes only the pole of this chart. Radial inversion preserves \(\theta\) and replaces \(r\) by \(r^{-1}\).

\begin{proposition}[Reflection about the line of midpoint real part]
\label{hrr:reflexion-compleja} In the preceding domain,
\begin{align}
 s(r^{-1},\theta)&=1-\overline{s(r,\theta)},
 \label{hrr:reflexion}\\
 \operatorname{Re}s(r,\theta)-\frac12
 &=\frac{1-r^2}{2|1-re^{i\theta}|^2}.
 \label{hrr:reflexion-parte-real}
\end{align}
Therefore, the circle \(r=1\), apart from the pole, maps to the line \(\operatorname{Re}s=1/2\). Radial inversion is realized as reflection about that vertical line.
\end{proposition}
\begin{proof}
The conjugate of \(s(r,\theta)\) is \((1-re^{-i\theta})^{-1}\). By direct calculation,
\[
 1-\frac1{1-re^{-i\theta}}
 =\frac{-re^{-i\theta}}{1-re^{-i\theta}}
 =\frac1{1-r^{-1}e^{i\theta}},
\]
proving \eqref{hrr:reflexion}. For the second identity, multiply numerator and denominator of \(s\) by \(1-re^{-i\theta}\). Its real part is
\[
 \frac{1-r\cos\theta}{1-2r\cos\theta+r^2}.
\]
Subtracting \(1/2\) reduces the numerator to \((1-r^2)/2\). The denominator is positive on the domain. It follows that \(\operatorname{Re}s=1/2\) is equivalent to \(r=1\). Finally, if \(s=x+iy\), then \(1-\overline s=(1-x)+iy\), which is exactly reflection about the stated vertical line.
\end{proof}

For \(\theta=\pi\), the radii \(r=2\) and \(r=1/2\) produce \(s=1/3\) and \(s=2/3\) respectively. The point \(r=1\) produces \(s=1/2\). This example makes reciprocity visible through three rational chart values.

The proved transformation belongs to the geometry of the complex chart. Applying it to a spectral operator requires preserving the operator and its domain as well as that chart. Likewise, the radial chart \(R_*\) preserves the angular section generating it. Equality of an inversion operation in two realizations supplies the stated correspondence; identification of their antecedent objects is maintained through the transports already constructed.

\section{Recovery of the flow and compatibility of readers}
\label{hrr:archivo}

Applying these readers to a flow direction requires preserving the generated direction as well as its scalar values. The reduced archive previously constructed has operators \(M\) and \(L\) with \(LM=P_{11}\). If \(m=MK\), then
\begin{equation}
 u_K=P_3Lm.
 \label{hrr:archivo-direccion}
\end{equation}
This equality preserves the action of each operator and allows recovery to be composed with the flow.

\begin{proposition}[Reconstruction of the flow from the reduced archive]
\label{hrr:archivo-flujo} With the preceding operators and marking, the archive \(m\) and a positive mark \(r\) determine
\[
 g_{K,r}(t)=
 \operatorname{diag}\left(
 r^{\,t(P_3Lm)_1/\|P_3Lm\|},\ldots,
 r^{\,t(P_3Lm)_{12}/\|P_3Lm\|}\right).
\]
All composition identities in \eqref{hrr:producto}--\eqref{hrr:determinante} are preserved.
\end{proposition}
\begin{proof}
From \(m=MK\) and \(LM=P_{11}\) follows \(P_3Lm=P_3LMK=P_3P_{11}K=u_K\). The direction is nonzero by the already proved property of the generated register. Normalization recovers \(v\) exactly, and its substitution into \eqref{hrr:g-definicion} recovers every flow entry. The identities depend only on this direction and positivity of marks, so remain valid.
\end{proof}

The archive identity recovers the direction of deformation. Regional readers supply their marks. The diagonal reader of the radial shape preserves its absolute anisotropy, and the oriented register recovers the sign. The resulting composition assigns invariance and reversibility to explicit data:
\[
 \begin{gathered}
 \text{sheet returns}
 \longrightarrow(\Theta_{\pi,n},\Theta_{e,n})
 \longrightarrow\pi/e,\\
 \text{archive of }K\longrightarrow u_K
 \longrightarrow g_{K,r}(t)
 \longrightarrow\text{covolume and transported minors},\\
 (s,D,\varepsilon)\longrightarrow(a,b)
 \longleftrightarrow(R_*,w,x_{\mathrm K}).
 \end{gathered}
\]
The first line distinguishes the common support of two marks. The second expresses the register's operator function. The third specifies the geometric information needed for inversion. The enriched state preserves phase, turns, carries and incidences allowing these readers to be composed within their genealogy.

\section{Temporal realization and domain of physical interpretation}

Nonadic holonomy preserves its action on phase and memory; phase return increments the turn register. Counts of the modal envelope and mark ratios are evaluated on subsequent realizations. Radial reciprocity acts on a shape chart and its orientation. The proved articulation combines these operations through compatible readers, preserving the distinction between their parameters.

The relation \(\pi/\varphi^2\) appears as a marked return ratio and as the scale of the conservative flow. The ratio \(\pi/e\) appears as a mark quotient on the same support. Both satisfy the exact factorization \eqref{hrr:composicion-limite}. Their involvement in a gravitational or cosmological realization preserves the conditions of the corresponding physical reader: clock identification, dimensional quantities, boundary response and geometric domain. This chapter's result concerns mathematical composition of the generated readers. Their physical applications subsequently use the corresponding realization.
